% =============================================================================
% Presentación Beamer P02 – Calibración del Túnel de Viento
% Técnicas de Medida – Ingeniería Aeronáutica
% =============================================================================
\documentclass[aspectratio=169,12pt]{beamer}

\usetheme{Madrid}
\usecolortheme{seahorse}
\usefonttheme{professionalfonts}

\usepackage[utf8]{inputenc}
\usepackage[T1]{fontenc}
\usepackage[spanish]{babel}
\usepackage{amsmath,amssymb}
\usepackage{siunitx}
\sisetup{output-decimal-marker={,}}
\usepackage{tikz}
\usetikzlibrary{arrows.meta,positioning,calc,shapes.geometric}
\usepackage{circuitikz}
\usepackage{booktabs}
\usepackage{listings}
\lstset{
  basicstyle=\ttfamily\tiny,
  keywordstyle=\color{blue!80}\bfseries,
  commentstyle=\color{green!50!black},
  backgroundcolor=\color{gray!10},
  frame=single, breaklines=true
}

\definecolor{aeroblue}{RGB}{0,51,102}
\definecolor{aerored}{RGB}{153,0,0}
\definecolor{aerogreen}{RGB}{0,102,51}

\setbeamercolor{frametitle}{bg=aeroblue,fg=white}
\setbeamercolor{title}{bg=aeroblue,fg=white}
\setbeamercolor{structure}{fg=aeroblue}

\title[P02 -- Calibración Túnel]{P02: Calibración del Túnel de Viento}
\subtitle{Técnicas de Medida -- Ingeniería Aeronáutica}
\author{Laboratorio de Técnicas de Medida}
\institute{Departamento de Ingeniería Aeroespacial}
\date{\today}

% =============================================================================
\begin{document}

\begin{frame}
  \titlepage
\end{frame}

\begin{frame}{Contenido}
  \tableofcontents
\end{frame}

% =============================================================================
\section{Motivación}

\begin{frame}{¿Por qué calibrar el túnel?}
  \begin{columns}
    \begin{column}{0.55\textwidth}
      \begin{itemize}
        \item Ensayos aerodinámicos requieren conocer $V_\infty$ con
              \textbf{incertidumbre $< 1\%$}
        \item El variador de frecuencia (VFD) no da velocidad directamente
        \item $Re$, $C_L$, $C_D$ dependen críticamente de $V_\infty$
      \end{itemize}
      \vspace{0.5em}
      \begin{alertblock}{Objetivo}
        Curva $V = f(f_\mathrm{VFD})$ con incertidumbre cuantificada.
      \end{alertblock}
      \vspace{0.5em}
      \textbf{Túnel del laboratorio:}\\
      Sección $\SI{30}{\cm}\times\SI{30}{\cm}$,
      $V_\mathrm{max}\approx\SI{50}{\m\per\s}$,\\
      VFD $f \in [5,55]\,\si{\Hz}$
    \end{column}
    \begin{column}{0.42\textwidth}
      \begin{tikzpicture}[scale=0.75,>=Stealth]
        \draw[aeroblue,very thick] (0,1.0) -- (6,1.0);
        \draw[aeroblue,very thick] (0,-1.0) -- (6,-1.0);
        \draw[fill=gray!15,draw=aeroblue] (0,-1.0) rectangle (1.0,1.0);
        \draw[fill=gray!20,draw=aeroblue] (1.0,1.0) -- (2.5,0.55)
              -- (2.5,-0.55) -- (1.0,-1.0) -- cycle;
        \draw[fill=gray!10,draw=aeroblue] (2.5,0.55) -- (4.5,0.55)
              -- (4.5,-0.55) -- (2.5,-0.55) -- cycle;
        \draw[fill=gray!20,draw=aeroblue] (4.5,0.55) -- (6.0,1.0)
              -- (6.0,-1.0) -- (4.5,-0.55) -- cycle;
        \foreach \y in {-0.3,0,0.3} {
          \draw[->,blue!60] (2.7,\y) -- (3.5,\y);
        }
        \draw[fill=gray!70,draw=black,very thick] (3.5,-0.55) -- (3.5,-0.05);
        \fill[aerored] (3.5,-0.05) circle (0.08);
        \node[above,font=\tiny] at (3.5,0) {Pitot};
        \node[below,font=\tiny,aeroblue] at (3.5,-0.65) {Secc.~prueba};
      \end{tikzpicture}
    \end{column}
  \end{columns}
\end{frame}

% =============================================================================
\section{Teoría: Tubo de Pitot}

\begin{frame}{Principio de Bernoulli -- Tubo de Pitot}
  \begin{columns}
    \begin{column}{0.5\textwidth}
      \textbf{Ecuación de Bernoulli (incompresible):}
      \begin{equation*}
        p + \tfrac{1}{2}\rho V^2 = \text{cte}
      \end{equation*}

      \textbf{En el punto de remanso:} $V = 0$
      \begin{equation*}
        p_0 = p_\infty + \tfrac{1}{2}\rho V_\infty^2
      \end{equation*}

      \vspace{0.3em}
      \begin{block}{Velocidad Pitot}
        \begin{equation*}
          \boxed{V_\infty = \sqrt{\frac{2\,\Delta p}{\rho}}}
        \end{equation*}
        $\Delta p = p_0 - p_s$, medida con manómetro
      \end{block}
    \end{column}
    \begin{column}{0.47\textwidth}
      \begin{tikzpicture}[scale=0.8,>=Stealth]
        % Flow arrows
        \foreach \y in {0.6,0.3,0,-0.3,-0.6} {
          \draw[->,blue!60,thick] (0,\y) -- (1.5,\y);
        }
        % Pitot body
        \draw[fill=gray!40,draw=black,thick]
          (1.5,-0.5) rectangle (1.8,0.5);
        % Stagnation point
        \fill[aerored] (1.8,0) circle (0.1);
        \node[aerored,right,font=\tiny] at (1.9,0) {$p_0$ (total)};
        % Static tap
        \draw[aerored!60,dashed,thick] (1.65,0.3) -- (2.5,0.3);
        \node[right,font=\tiny] at (2.5,0.3) {$p_s$ (estática)};
        % Differential
        \draw[<->,gray] (2.1,-0.4) -- (2.1,0.4)
              node[midway,right,font=\tiny] {$\Delta p$};
        \node[left,font=\tiny,blue!60] at (0.7,0.6) {$V_\infty$};
        % Stagnation streamline
        \draw[aerored!60,thick,->] (0.5,0) -- (1.75,0);
      \end{tikzpicture}
    \end{column}
  \end{columns}
\end{frame}

\begin{frame}{Densidad del aire y compresibilidad}
  \begin{columns}
    \begin{column}{0.5\textwidth}
      \textbf{Densidad (gas ideal):}
      \begin{equation*}
        \rho = \frac{p_\mathrm{atm}}{R_\mathrm{aire}\,T},
        \quad R_\mathrm{aire} = \SI{287}{\J\per\kg\per\K}
      \end{equation*}

      \vspace{0.3em}
      \textbf{Corrección por compresibilidad:}
      Para $Ma < 0.3$ ($V < \SI{100}{\m\per\s}$):
      \begin{equation*}
        \varepsilon = \frac{V - V_\mathrm{inc}}{V_\mathrm{inc}} < 1\%
      \end{equation*}
      \textit{En nuestro túnel: $\varepsilon < 0.3\%$. Se puede omitir.}
    \end{column}
    \begin{column}{0.47\textwidth}
      \begin{tikzpicture}[scale=0.85]
        \draw[->] (0,0) -- (5,0) node[right,font=\tiny] {$Ma$};
        \draw[->] (0,0) -- (0,3) node[above,font=\tiny] {Error (\%)};
        \draw[aeroblue,thick] (0,0) .. controls (1,0.02) and (2.5,0.1)
              .. (3,0.5) .. controls (3.5,1.0) .. (4.5,2.5);
        \draw[aerored,dashed] (1.5,0) -- (1.5,3);
        \node[aerored,below,font=\tiny] at (1.5,0) {$Ma=0.3$};
        \draw[gray,dashed] (0,0.15) -- (1.5,0.15);
        \node[gray,left,font=\tiny] at (0,0.15) {1\%};
        \fill[aeroblue!20,opacity=0.5] (0,0) rectangle (1.5,3);
        \node[aeroblue,font=\tiny] at (0.75,2.5) {Nuestro\\rango};
        \draw[->,aerored] (0.5,0.1) node[below,font=\tiny] {$Ma_\mathrm{max}=0.15$}
              (0.5,0) -- (0.5,0.1);
      \end{tikzpicture}
    \end{column}
  \end{columns}
\end{frame}

% =============================================================================
\section{Cálculo Numérico}

\begin{frame}{Calibración numérica -- Ejemplo}
  \begin{columns}
    \begin{column}{0.5\textwidth}
      \textbf{Datos ($\rho = \SI{1.225}{\kg\per\m\cubed}$):}
      \begin{table}
        \centering
        \tiny
        \begin{tabular}{SSS}
          \toprule
          {$f$ (Hz)} & {$\Delta p$ (Pa)} & {$V$ (m/s)} \\
          \midrule
           5 &   16.6 &  5.2 \\
          15 &  149.2 & 15.6 \\
          25 &  414.4 & 26.0 \\
          35 &  811.9 & 36.4 \\
          45 & 1342   & 46.8 \\
          55 & 2004   & 57.2 \\
          \bottomrule
        \end{tabular}
      \end{table}
      \vspace{0.3em}
      \textbf{Ajuste lineal:}
      \begin{equation*}
        V(f) = 1.044\,f - 0.002 \quad R^2 = 0.9997
      \end{equation*}
    \end{column}
    \begin{column}{0.47\textwidth}
      \begin{tikzpicture}[scale=0.75]
        \draw[->] (0,0) -- (6,0) node[right,font=\tiny] {$f$ (Hz)};
        \draw[->] (0,0) -- (0,5) node[above,font=\tiny] {$V$ (m/s)};
        % Grid
        \foreach \x in {1,2,3,4,5} {
          \draw[gray!30] (\x,0) -- (\x,5);
          \node[below,font=\tiny] at (\x,0) {\pgfmathprintnumber[precision=0]{\x0}};
        }
        \foreach \y in {1,2,3,4} {
          \draw[gray!30] (0,\y) -- (6,\y);
          \node[left,font=\tiny] at (0,\y) {\pgfmathprintnumber[precision=0]{\y5}};
        }
        % Data points
        \fill[aerored] (0.5,1.04) circle (0.08);
        \fill[aerored] (1.5,3.12) circle (0.08);
        \fill[aerored] (2.5,5.2/12) circle (0.08);
        % Trend line
        \draw[aeroblue,thick] (0,0) -- (5.5,4.78);
        \node[aeroblue,font=\tiny,rotate=42] at (3.5,3.0) {$V=1.044f$};
        % Labels y
        \node[left,font=\tiny] at (0,1.04) {5.2};
        \node[left,font=\tiny] at (0,2.34) {26.0};
        \node[left,font=\tiny] at (0,4.34) {46.8};
        % Points all
        \foreach \f/\v in {5/5.2, 10/10.4, 15/15.6, 20/20.8, 25/26.0, 35/36.4, 45/46.8, 55/57.2} {
          \pgfmathsetmacro{\fx}{\f/10}
          \pgfmathsetmacro{\vy}{\v/12}
          \fill[aerored] (\fx,\vy) circle (0.07);
        }
      \end{tikzpicture}
    \end{column}
  \end{columns}
\end{frame}

\begin{frame}{Incertidumbre de velocidad}
  \textbf{Propagación de errores (GUM):}
  \begin{equation*}
    \frac{u_V}{V} = \frac{1}{2}\sqrt{
      \left(\frac{u_{\Delta p}}{\Delta p}\right)^2
      + \left(\frac{u_\rho}{\rho}\right)^2}
  \end{equation*}

  \vspace{0.3em}
  \textbf{Ejemplo: $V = \SI{28.57}{\m\per\s}$ ($\Delta p = \SI{500}{\Pa}$)}

  \begin{columns}
    \begin{column}{0.5\textwidth}
      \begin{itemize}
        \item $u_{\Delta p} = \SI{2}{\Pa}$ (manómetro)
        \item $u_\rho/\rho = 0.3\%$ ($\pm\SI{1}{\K}$)
      \end{itemize}
      \begin{align*}
        \frac{u_V}{V} &= \frac{1}{2}\sqrt{(0.004)^2 + (0.003)^2}\\
                      &\approx 0.0025 = 0.25\%
      \end{align*}
    \end{column}
    \begin{column}{0.47\textwidth}
      \begin{alertblock}{Resultado}
        \centering
        $V = (28.57 \pm 0.07)\,\si{\m\per\s}$\\
        (95\% confianza, $k=2$)
      \end{alertblock}
      \vspace{0.3em}
      Adecuado para:
      \begin{itemize}\small
        \item Ensayos polares de perfil alar
        \item Medición de $Re$ con $<1\%$ de error
        \item Calibración de HWA y LDA
      \end{itemize}
    \end{column}
  \end{columns}
\end{frame}

% =============================================================================
\section{Caso de Estudio}

\begin{frame}{Caso: Polar aerodinámica de NACA 2412}
  \begin{block}{Escenario}
    Calibración para ensayo de perfil NACA~2412 ($c = \SI{100}{\mm}$)
    a $Re_c = 3\times10^5$.
  \end{block}
  \begin{columns}
    \begin{column}{0.5\textwidth}
      \textbf{Velocidad requerida:}
      \begin{equation*}
        V = \frac{Re_c\,\nu}{c} = \frac{3\!\times\!10^5 \times 1.5\!\times\!10^{-5}}{0.1}
          = \SI{45}{\m\per\s}
      \end{equation*}
      \textbf{Frecuencia del VFD:}
      \begin{equation*}
        f = \frac{45.0}{1.044} \approx \SI{43.1}{\Hz}
      \end{equation*}
      \textbf{$\Delta p$ esperada:}
      \begin{equation*}
        \Delta p = \frac{1.225\times45^2}{2} \approx \SI{1240}{\Pa}
      \end{equation*}
    \end{column}
    \begin{column}{0.47\textwidth}
      \begin{tikzpicture}[scale=0.8]
        % Airfoil NACA 2412 approximate shape
        \draw[aeroblue,very thick]
          (0,0) .. controls (0.5,0.35) and (1.0,0.4) .. (1.5,0.25)
               .. controls (2.0,0.1) .. (3.0,0)
               .. controls (2.5,-0.05) and (1.0,-0.08) .. (0,0);
        \draw[->,gray!60,thick] (-0.5,0.15) -- (0.5,0.15)
              node[right,font=\tiny] {$U_\infty=\SI{45}{\m\per\s}$};
        \draw[<->,gray] (0,-0.25) -- (3.0,-0.25)
              node[midway,below,font=\tiny] {$c=\SI{100}{\mm}$};
        \node[aeroblue,above,font=\small] at (1.5,0.5) {NACA 2412};
        \node[below,font=\tiny] at (1.5,-0.4) {$Re_c = 3\times10^5$};
      \end{tikzpicture}
    \end{column}
  \end{columns}
\end{frame}

% =============================================================================
\section{Notebook Jupyter}

\begin{frame}[fragile]{P02\_Tunel\_Calibracion.ipynb}
  \begin{columns}
    \begin{column}{0.55\textwidth}
      \textbf{Ruta:}\\
      \texttt{\small Practicas/P02\_Calibracion\_Tunel/\\notebooks/P02\_Tunel\_Calibracion.ipynb}

      \vspace{0.5em}
      \textbf{Funciones del módulo \texttt{calib\_utils.py}:}
      \begin{itemize}\small
        \item \texttt{velocity\_from\_dp(dp, rho)} -- ec. Bernoulli
        \item Lectura automática de \texttt{files/p2/}
        \item Regresión polinomial con \texttt{numpy.polyfit}
        \item Propagación de incertidumbre (GUM)
        \item Exportación a CSV y gráficas
      \end{itemize}
    \end{column}
    \begin{column}{0.42\textwidth}
      \begin{lstlisting}[language=Python]
from calib_utils import velocity_from_dp
import numpy as np

dp  = np.array([16.6, 149.2,
                414.4, 1342, 2004])
rho = 1.225  # kg/m^3

V = velocity_from_dp(dp, rho)
# array([ 5.2, 15.6, 26.0,
#         46.8, 57.2])

f = np.array([5, 15, 25, 45, 55])
c = np.polyfit(f, V, deg=1)
print(f"V = {c[0]:.3f}*f + {c[1]:.3f}")
# V = 1.044*f + (-0.002)
      \end{lstlisting}
    \end{column}
  \end{columns}
\end{frame}

% =============================================================================
\section{Resumen}

\begin{frame}{Conclusiones}
  \begin{enumerate}
    \item La calibración establece $V = 1.044\,f_\mathrm{VFD}$
          con $R^2 = 0.9997$ (relación \textbf{lineal}).
    \item La incertidumbre relativa $u_V/V < 0.3\%$, adecuada para
          ensayos aeronáuticos de precisión.
    \item La \textbf{uniformidad} del perfil $I_u < 1\%$ en el núcleo
          central confirma la calidad aerodinámica del túnel.
    \item El módulo \texttt{calib\_utils.py} y el notebook automatizan
          todo el proceso de calibración.
  \end{enumerate}
  \vspace{0.5em}
  \begin{block}{Recursos}
    Manual: \texttt{manuals/P02\_Calibracion\_Tunel/manual\_P02.tex}\\
    Notebook: \texttt{Practicas/P02\_Calibracion\_Tunel/notebooks/P02\_Tunel\_Calibracion.ipynb}\\
    Código: \texttt{Practicas/P02\_Calibracion\_Tunel/src/calib\_utils.py}
  \end{block}
\end{frame}

\begin{frame}{Referencias}
  \begin{itemize}
    \item Barlow, Rae \& Pope (1999). \textit{Low-Speed Wind Tunnel Testing}, 3ª ed., Wiley
    \item Anderson, J.D. (2017). \textit{Fundamentals of Aerodynamics}, 6ª ed., McGraw-Hill
    \item JCGM 100:2008 -- \textit{Guide to the Expression of Uncertainty in Measurement}
  \end{itemize}
\end{frame}

\end{document}
